\documentclass[a4paper,12pt]{article}
\usepackage{graphicx}
\usepackage{caption}
\usepackage{subcaption}
\usepackage{float}
\usepackage[margin=1in]{geometry}
\usepackage{lmodern}
\usepackage{amsmath}
\usepackage{xcolor}
\usepackage{listings}
\usepackage{pgffor}

\graphicspath{{../../python/bayes/breastcancer/output/}}

\begin{document}

\title{Relatório Bayes}
\author{Eduardo}
\date{\today}
\maketitle

\section*{Covariancia}

\begin{figure}[H]
  \centering
  \includegraphics[width=0.75\textwidth]{default_likelihood.png}
  \caption{Espaço de verossimilhança para o modelo padrão}
  \label{fig:likelihood_def}
\end{figure}

\section*{Diagonal}

Assume que todas as variáveis são independentes e usa o mesmo raio em todos os eixos.

\begin{figure}[H]
  \centering
  \includegraphics[width=0.75\textwidth]{accuracy_diagonal.png}
  \caption{Precisão vs raio}
  \label{fig:accuracy_diag}
\end{figure}

\foreach \case in {best, mean, worst} {
  \begin{figure}[H]
    \centering
    \includegraphics[width=0.75\textwidth]{likelihood_diagonal_\case.png}
    \caption{Espaço de verossimilhança para o modelo diagonal (\case)}
    \label{fig:likelihood_diag_\case}
  \end{figure}
}

\section*{Completa}

Monta a matriz de convariância usando todo o conjunto de dados, em oposição à padrão, que considera os dados de cada classe separadamente.

\begin{figure}[H]
  \centering
  \includegraphics[width=0.75\textwidth]{accuracy_full.png}
  \caption{Precisão vs raio}
  \label{fig:accuracy_full}
\end{figure}

\foreach \case in {best, mean, worst} {
  \begin{figure}[H]
    \centering
    \includegraphics[width=0.75\textwidth]{likelihood_full_\case.png}
    \caption{Espaço de verossimilhança para o modelo completo (\case)}
    \label{fig:likelihood_full_\case}
  \end{figure}
}

\section*{Diagonal com variança por atributo}

Considera independência entre os atributos, mas preserva a variança de cada um.

\begin{figure}[H]
  \centering
  \includegraphics[width=0.75\textwidth]{accuracy_per_feature.png}
  \caption{Precisão vs raio}
  \label{fig:accuracy_per_feature}
\end{figure}

\foreach \case in {best, mean, worst} {
  \begin{figure}[H]
    \centering
    \includegraphics[width=0.75\textwidth]{likelihood_per_feature_\case.png}
    \caption{Espaço de verossimilhança para o modelo diagonal com variança por atributo (\case)}
    \label{fig:likelihood_per_feature_\case}
  \end{figure}
}

\section*{Completa Regularizada}

Monta a matriz de convariância usando todo o conjunto de dados e adiciona um fator de regularização.

\begin{figure}[H]
  \centering
  \includegraphics[width=0.75\textwidth]{accuracy_regularized.png}
  \caption{Precisão vs raio}
  \label{fig:accuracy_regularized}
\end{figure}

\foreach \case in {best, mean, worst} {
  \begin{figure}[H]
    \centering
    \includegraphics[width=0.75\textwidth]{likelihood_regularized_\case.png}
    \caption{Espaço de verossimilhança para o modelo completo regularizado (\case)}
    \label{fig:likelihood_regularized_\case}
  \end{figure}
}

\end{document}
